\chapter{Convolución periódica y Fourier finita}\label{cap:v4-05}
La convolución combina una función con todas las traslaciones de otra. En el círculo sirve para suavizar y aproximar; en un grupo finito es una suma que se puede calcular sin cuestiones de convergencia. Ambos modelos muestran la misma identidad: Fourier convierte convolución en producto.

\section{El álgebra de convolución del círculo}\label{lectura:v4-08}
\index{convolucion@convolución!periodica@periódica}
Para $f,g\in L^1(\T)$ definimos, donde la integral sea absolutamente convergente,
\[
 (f*g)(x)=\int_\T f(t)g(x-t)\,dm(t).
\]
La lectura 8 utiliza $dt$ y cambia su norma $L^1$ por la no normalizada. Aquí mantenemos $dm=dt/(2\pi)$ en todas las normas y convoluciones periódicas; las dos convenciones se relacionan multiplicando la convolución por $2\pi$.
\begin{teorema}\label{v4:conv-T}
La convolución está definida casi en todas partes, pertenece a $L^1(\T)$ y satisface
\[
 \norm{f*g}_1\leq\norm f_1\norm g_1,\qquad
 \int_\T f*g\,dm=\left(\int_\T f\,dm\right)\left(\int_\T g\,dm\right).
\]
Además $\widehat{f*g}(n)=\widehat f(n)\widehat g(n)$.
\end{teorema}
\begin{proof}
Elijamos representantes borelianos de $f,g$, modificándolos en conjuntos nulos si hace falta. El integrando $(x,t)\mapsto f(t)g(x-t)$ es medible. Las modificaciones nulas no afectan al resultado porque las traslaciones preservan nulos y el cálculo de Tonelli siguiente controla las secciones excepcionales. Por invariancia de $dm$,
\[
 \int_\T\int_\T |f(t)g(x-t)|\,dm(x)dm(t)
 =\int_\T|f(t)|\norm g_1\,dm(t)=\norm f_1\norm g_1.
\]
Fubini da existencia casi en todas partes y medibilidad de la integral. La desigualdad triangular, seguida de integración en $x$, da la cota. Aplicar Fubini sin valores absolutos da la fórmula para la integral. Finalmente, el factor $e^{-inx}$ tiene módulo uno y no altera la integrabilidad absoluta; con $u=x-t$,
\[
 \widehat{f*g}(n)=\int_\T\int_\T f(t)g(u)e^{-in(t+u)}\,dm(u)dm(t)
 =\widehat f(n)\widehat g(n).
\]
No se identifica la $\sigma$-álgebra producto de dos completaciones con toda la completación del producto: el uso de representantes borelianos evita esa afirmación incorrecta de la lectura.
\end{proof}

\section{Identidades aproximadas y unicidad}
\begin{lema}\label{v4:traslacion-T}
Si $1\leq p<\infty$ y $f\in L^p(\T)$, entonces $\norm{f(\cdot-t)-f}_p\to0$ cuando $t\to0$. La aplicación de traslación es uniformemente continua como aplicación de $\T$ en $L^p(\T)$.
\end{lema}
\begin{proof}
Para una función continua la conclusión se deduce de su continuidad uniforme y $\norm h_p\leq\norm h_\infty$. Para $f$ general escojamos continua $g$ con $\norm{f-g}_p<\eps/3$. La invariancia de la norma por traslación da
\[
 \norm{f(\cdot-t)-f}_p\leq2\norm{f-g}_p+\norm{g(\cdot-t)-g}_p.
\]
El último término es pequeño para $t$ pequeño. La diferencia entre dos traslaciones solo depende de la diferencia de sus parámetros, lo que da uniformidad.
\end{proof}
\begin{definicion}
Una sucesión de Dirac en $\T$ es una sucesión $k_j\geq0$ de $L^1$ con integral uno y $\int_{|t|\geq\delta}k_j\,dm\to0$ para cada $0<\delta<\pi$. Es fuerte si cada $k_j$ es esencialmente acotada y $\sup_{\delta\leq|t|\leq\pi}k_j(t)\to0$.
\end{definicion}
\begin{ejemplo}
Los núcleos $K_N$ de Fejér forman una sucesión fuerte. Otro ejemplo es
\[
 k_j(t)=\frac{\cos^{2j}(t/2)}{\int_\T\cos^{2j}(u/2)\,dm(u)}.
\]
Para $|t|\geq\delta$, el numerador es a lo sumo $\cos^{2j}(\delta/2)$. Al integrar el denominador sólo sobre $[-\delta/2,\delta/2]$, obtenemos
\[
 k_j(t)\leq\frac1{m([-\delta/2,\delta/2])}
 \left(\frac{\cos(\delta/2)}{\cos(\delta/4)}\right)^{2j}.
\]
El cociente de cosenos es menor que uno, así que la cota tiende a cero. Positividad e integral uno son inmediatas; la cota uniforme de la cola implica concentración de masa.
\end{ejemplo}
\begin{teorema}\label{v4:dirac-T}
Si $(k_j)$ es una sucesión de Dirac, entonces $f*k_j\to f$ en $L^p(\T)$ para $1\leq p<\infty$. Para $f$ continua, la convergencia es uniforme.
\end{teorema}
\begin{proof}
Para $p=1$, Tonelli y la desigualdad triangular dan
\[
 \norm{f*k_j-f}_1\leq\int_\T k_j(t)\norm{f(\cdot-t)-f}_1\,dm(t).
\]
Para $1<p<\infty$, Jensen aplicado a la probabilidad $k_jdm$ da
\[
 \norm{f*k_j-f}_p^p\leq\int_\T k_j(t)\norm{f(\cdot-t)-f}_p^p\,dm(t).
\]
En $|t|<\delta$ la diferencia de normas es uniformemente pequeña por el lema anterior; fuera, su potencia está acotada por $(2\norm f_p)^p$ y la masa tiende a cero. Esto prueba convergencia en norma. Para $f$ continua se aplica el mismo argumento al supremo de $|f(x-t)-f(x)|$, pequeño por continuidad uniforme.
\end{proof}
\begin{corolario}\label{v4:unicidad-T}
Si $f\in L^1(\T)$ tiene todos sus coeficientes nulos, entonces $f=0$ casi en todas partes. El álgebra de convolución $L^1(\T)$ no tiene unidad.
\end{corolario}
\begin{proof}
La fórmula finita de Fejér muestra que $f*K_N=0$ cuando todos los coeficientes de $f$ son cero. El teorema anterior da $\norm f_1=0$. Si hubiera unidad $u$, convolucionar con $e_n$ daría $\widehat u(n)e_n=e_n$ para todo $n$, por lo que $\widehat u(n)=1$. Esto contradice Riemann--Lebesgue.
\end{proof}
\begin{proposicion}
La aplicación de coeficientes $L^1(\T)\to c_0(\Z)$ no es sobreyectiva.
\end{proposicion}
\begin{proof}
Si fuera sobreyectiva, sería un isomorfismo lineal continuo entre espacios de Banach por el teorema de la inversa acotada, pues acabamos de probar inyectividad. Existiría $C$ tal que $\norm f_1\leq C\sup_n|\widehat f(n)|$. Para $f=D_N$, el supremo de los coeficientes es uno, pero $\norm{D_N}_1\to\infty$. Verifiquemos esto último. Con $M=N+1/2$ y $\sin(t/2)\leq t/2$,
\[
 \norm{D_N}_1\geq\frac2\pi\int_0^1\frac{|\sin(Mt)|}{t}\,dt.
\]
En cada intervalo $[k\pi/M,(k+1)\pi/M]\subset[0,1]$, $k\geq1$, la integral es al menos
\[
 \frac{M}{(k+1)\pi}\int_{k\pi/M}^{(k+1)\pi/M}|\sin(Mt)|\,dt
 =\frac2{(k+1)\pi}.
\]
El número de esos intervalos crece linealmente con $N$ y la suma armónica diverge. Esto contradice la supuesta cota de la inversa y completa la afirmación anunciada en la lectura 7.
\end{proof}

\section{Ejercicios de convolución periódica}
\begin{ejercicio}[Lectura 8, ejercicio 1]\label{ej:v4-l08-1}
Demostrar que $L^1(\T)$, con convolución, es un álgebra de Banach conmutativa.
\end{ejercicio}
\begin{solucion}
La bilinealidad sigue de la linealidad integral y la cota submultiplicativa del teorema~\ref{v4:conv-T}. La completitud es la de $L^1$. Cambiar $t$ por $x-t$ prueba conmutatividad. Para asociatividad, la integral triple del valor absoluto está acotada por $\norm f_1\norm g_1\norm h_1$. Fubini permite entonces calcular casi en todas partes
\[
 ((f*g)*h)(x)=\int_\T\int_\T f(s)g(t-s)h(x-t)\,dm(s)dm(t)
 =\int_\T f(s)\int_\T g(u)h(x-s-u)\,dm(u)dm(s),
\]
con $u=t-s$. El último miembro es $f*(g*h)(x)$.
\end{solucion}
\begin{ejercicio}[Lectura 8, ejercicio 2: cota del operador]\label{ej:v4-l08-2}
Si $f\in L^1(\T)$, $g\in L^p(\T)$ y $1<p<\infty$, demostrar que $f*g\in L^p$ y $\norm{f*g}_p\leq\norm f_1\norm g_p$. Precisar la norma de $T_f:g\mapsto f*g$.
\end{ejercicio}
\begin{solucion}
Si $\norm f_1=0$ la conclusión es inmediata. En otro caso, Hölder para la medida finita $|f(t)|dm(t)$ da
\[
 |f*g(x)|^p\leq\norm f_1^{p-1}\int_\T |f(t)|\,|g(x-t)|^pdm(t).
\]
Tonelli e invariancia de la medida hacen integrable el lado derecho y, tras integrar, dan $\norm{f*g}_p^p\leq\norm f_1^p\norm g_p^p$. En particular la integral de convolución existe absolutamente casi en todas partes. Por tanto $\norm{T_f}\leq\norm f_1$.

La igualdad afirmada en el original es falsa en general. Para $p=2$ y $f(t)=1+e^{it}$, Fourier muestra que $T_f$ es la proyección ortogonal sobre $\operatorname{span}\{1,e^{it}\}$, de norma uno, mientras que
\[
 \norm f_1=\frac1{2\pi}\int_0^{2\pi}2|\cos(t/2)|\,dt=\frac4\pi>1.
\]
La cota es la conclusión válida para todo $p$.
\end{solucion}
\begin{ejercicio}[Lectura 8, ejercicio 3]\label{ej:v4-l08-3}
Probar que $f\mapsto T_f$ es un homomorfismo continuo de álgebras de $L^1(\T)$ en $\Bcal(L^p(\T))$.
\end{ejercicio}
\begin{solucion}
La linealidad es inmediata. Si $f,h\in L^1$ y $g\in L^p$, aproximamos $g$ en $L^p$ por funciones acotadas $g_j$ (truncación en altura). Para cada $g_j$, las integrales de valor absoluto en la identidad de asociatividad son finitas porque el círculo tiene medida uno y $g_j$ es acotada. Así $T_{f*h}g_j=T_fT_hg_j$. La cota del ejercicio anterior permite pasar al límite en $L^p$ y obtener $T_{f*h}=T_fT_h$. También $\norm{T_f-T_h}\leq\norm{f-h}_1$, lo que prueba continuidad.
\end{solucion}
\begin{ejercicio}[Lectura 8, ejercicio 4]\label{ej:v4-l08-4}
Si $1<p<\infty$, $(k_j)$ es una sucesión de Dirac y $g\in L^p(\T)$, probar $\norm{g-g*k_j}_p\to0$.
\end{ejercicio}
\begin{solucion}
Jensen para la medida de probabilidad $k_jdm$ produce
\[
\norm{g-g*k_j}_p^p\leq\int k_j(t)\norm{g-g(\cdot-t)}_p^pdm(t).
\]
Para $\eps>0$ escojamos $\delta$ con $\norm{g-g(\cdot-t)}_p<\eps$ si $|t|<\delta$. La integral es a lo sumo
\[
 \eps^p+(2\norm g_p)^p\int_{|t|\geq\delta}k_j(t)\,dm(t).
\]
El límite superior es $\eps^p$; hacer $\eps\downarrow0$ da la conclusión. Esta es la versión $L^p$ completa de la sugerencia original, con una misma letra $g$ en hipótesis y conclusión.
\end{solucion}

\section{Un laboratorio finito: el grupo cíclico}
\index{Fourier!transformada discreta}\index{grupo!ciclico finito@cíclico finito}\index{matriz!circulante}
El complemento siguiente desarrolla el modelo finito presentado por Maximenko~\cite{maximenkoConvolucionFinita}. Sea $G=\Z/q\Z$, representado por $0,\ldots,q-1$, con medida de conteo. Los índices de una función $a:G\to\C$ se entienden módulo $q$. Definimos
\[
 (a*b)(r)=\sum_{k=0}^{q-1}a(k)b(r-k),\qquad
 \widehat a(j)=\sum_{k=0}^{q-1}a(k)e^{2\pi ijk/q}.
\]
\begin{teorema}\label{v4:dft}
Se cumplen inversión, Plancherel y el teorema de convolución:
\[
 a(k)=\frac1q\sum_{j=0}^{q-1}\widehat a(j)e^{-2\pi ijk/q},\qquad
 \sum_k|a(k)|^2=\frac1q\sum_j|\widehat a(j)|^2,
 \qquad\widehat{a*b}=\widehat a\,\widehat b.
\]
\end{teorema}
\begin{proof}
La suma geométrica da $\sum_{j=0}^{q-1}e^{2\pi ijr/q}=q$ si $r=0$ en $G$, y cero en otro caso. Sustituir la definición de $\widehat a$ en la fórmula inversa deja únicamente el término con índice $k$. Expandiendo $|\widehat a(j)|^2$ y sumando en $j$, la misma ortogonalidad elimina los términos cruzados y da $q\sum_k|a(k)|^2$.

Para convolución cambiamos $r=k+s$; para cada $k$, la aplicación $s\mapsto k+s$ permuta exactamente todos los elementos de $G$. Por tanto
\[
 \widehat{a*b}(j)=\sum_k\sum_s a(k)b(s)e^{2\pi ij(k+s)/q}
 =\widehat a(j)\widehat b(j).
\]
Esta biyección de índices es toda la justificación del cambio de suma; no hay ningún paso analítico oculto. El operador unitario $\mathcal F_q=q^{-1/2}\widehat{\phantom a}$ satisface $\mathcal F_q(a*b)=\sqrt q(\mathcal F_qa)(\mathcal F_qb)$: la normalización unitaria introduce el factor $\sqrt q$.
\end{proof}
La matriz de $T_a:b\mapsto a*b$ tiene entrada $a(r-k)$; cada columna es una traslación de la anterior. Se llama circulante. Los vectores $k\mapsto q^{-1/2}e^{-2\pi ijk/q}$ forman una base ortonormal de eigenvectores de $T_a$, con eigenvalores $\widehat a(j)$: sustituirlos en la definición de convolución lo verifica directamente. Así $\norm{T_a}_{\ell^2\to\ell^2}=\max_j|\widehat a(j)|$.
\begin{ejercicio}[Complemento: cálculo en cuatro puntos]\label{ej:v4-ampl-dft4}
En $\Z/4\Z$, calcular la convolución de $a=(1,2,0,0)$ y $b=(1,0,1,0)$ mediante sumas y mediante Fourier.
\end{ejercicio}
\begin{solucion}
Como $b$ vale uno en $0$ y $2$, $(a*b)(r)=a(r)+a(r-2)$, por lo que $a*b=(1,2,1,2)$. Las transformadas son
\[
 \widehat a=(3,1+2i,-1,1-2i),\qquad \widehat b=(2,0,2,0).
\]
El producto es $(6,0,-2,0)$. Invirtiendo, $(a*b)(k)=\tfrac14(6-2e^{-\pi ik})$, que vale uno si $k$ es par y dos si es impar. El cálculo exhibe la diferencia entre convolución cíclica y producto de polinomios sin reducción: aquí se trabaja módulo $z^4-1$.
\end{solucion}

\paragraph{La misma operación en tres modelos.}
Las sumas finitas, las series en grupos discretos y las integrales periódicas comparten la misma identidad de convolución. Al pasar a la recta desaparece la medida total finita; las hipótesis de integrabilidad y las estimaciones de Young serán las que aseguren que las operaciones sigan teniendo sentido.
